\section{Self-Attention：从两条序列对齐到序列内部关系}

传统 encoder--decoder attention 通常让 target-side query 去读取 source-side keys 与 values；self-attention 则令 query、key、value 来自同一个表示集合。这个变化看似只是在函数调用里把三个输入设成相同张量，统计含义却发生了改变：模型不再只做跨序列对齐，而开始学习一个样本内部的关系图。

\lecturefigure{slide-03-self-attention.jpg}{Self-attention 把 source 与 target 设为同一序列，学习 intra-sequence relations。}{Stanford Online 官方视频 00:01:42--00:02:39。}

设输入表示为 $X\in\mathbb{R}^{m\times d_{\mathrm{model}}}$，scaled dot-product attention 写为
\[
Q=XW_Q,\qquad K=XW_K,\qquad V=XW_V,
\]
\[
A=\operatorname{softmax}\!\left(\frac{QK^{\top}}{\sqrt{d_k}}\right),
\qquad Z=AV.
\]

\begin{itemize}
  \item $m$：序列长度，也就是参与关系建模的 token 数量；
  \item $W_Q,W_K,W_V$：把输入投影为 query、key、value 的可学习矩阵；
  \item $d_k$：每个 key/query 的维度，$\sqrt{d_k}$ 用于控制 dot product 的尺度；
  \item $A_{ij}$：第 $i$ 个 token 从第 $j$ 个 token 读取信息的归一化权重；
  \item $Z_i$：第 $i$ 个 token 聚合所有 value 后得到的新表示。
\end{itemize}

\teachervoice{讲者用 “blue ball” 说明 self-attention 的语义：模型应学会把形容词 blue 与它修饰的名词 ball 关联起来。这个例子不是在宣称某个 head 必然对应语法树，而是在强调 attention weight 可以承载输入内部的关系假设。}

\begin{warningbox}{Attention weight 不是自动获得的解释}
$A_{ij}$ 是一次前向计算中的信息混合权重。它能提示模型正在读取哪里，但不能单独证明某条语言学关系、因果关系或决策理由。解释时还需结合 value 内容、后续层、ablation 或 intervention。
\end{warningbox}

\subsection{Multi-Head Attention：并行学习不同关系子空间}

单个 attention map 只能在一个投影空间里比较 query 与 key。Multi-head attention 将表示投影到多个较低维子空间，让不同 heads 可以形成不同相似度度量，再把结果拼接并重新混合。读图时应先看“拆分--独立 attention--拼接”这条数据流，而不是把 head 数量理解为简单复制模型。

\lecturefigure{slide-04-multi-head-attention.jpg}{Multi-head attention：多个关系子空间独立计算后再拼接。}{Stanford Online 官方视频 00:02:40--00:03:27。}

\[
\operatorname{head}_r
=\operatorname{Att}\!\left(QW_r^Q,KW_r^K,VW_r^V\right),
\]
\[
\operatorname{MHA}(Q,K,V)
=\operatorname{Concat}(\operatorname{head}_1,\ldots,\operatorname{head}_H)W^O.
\]

\begin{itemize}
  \item $H$：attention head 数量；
  \item $W_r^Q,W_r^K,W_r^V$：第 $r$ 个 head 的独立投影；
  \item $W^O$：把所有 heads 的输出重新混合回 model dimension；
  \item 每个 head 可以学习不同关系，但没有机制保证人工可命名的完美分工。
\end{itemize}

\teachervoice{Gomez 的直觉是：一个 head 可以偏向形容词--名词关系，另一个 head 学别的模式。应把它读成“容量允许关系分工”，而不是“训练后每个 head 都会自动变成清晰的人类概念”。}

\begin{knowledgebox}{为什么除以 $\sqrt{d_k}$}
若 query 与 key 的各维近似独立且方差为常数，dot product 的方差随 $d_k$ 增大。过大的 logits 会让 softmax 过早饱和、梯度变小；缩放因子把数值范围拉回更稳定的区域。
\end{knowledgebox}

\subsection{Transformer 不是单一 trick}

Self-attention 和 multi-head attention 是容易被记住的标题，但 Gomez 特别强调原始系统还依赖一组当时并不显然的训练选择：LayerNorm、residual path、learning-rate warmup、初始化与 Tensor2Tensor 的工程默认值。论文发表后它们逐渐成为常识，容易让后来的读者误以为这些组件天然就该这样组合。

\teachervoice{团队曾反复做 ablation，发现移除 warmup、LayerNorm 或改动看似细小的配置会真实破坏优化。工程结论是：一个架构的成功往往来自彼此配合的机制集合，不能把所有贡献压缩为一张 attention 公式。}

\teachervoice{Gomez 回忆原始工作在约三个月内快速汇聚，并在截止日前冲刺完成；他当时并未预见影响。这段历史更有价值的读法不是制造“天才瞬间”，而是看到成熟工具、多人迭代、实验反馈与时机共同促成新范式。}

\subsection{本章小结}

Self-attention 把关系建模从两条序列之间移到同一输入内部；multi-head attention 提供多个关系子空间。Transformer 的实际可训练性还依赖一套优化与残差工程，不能只由 $QK^{\top}V$ 一式解释。

